\section{描述统计}

拿到一批数，先别急着推断，得把它讲清楚。这一章全是在干这件事：列成表、画成图、再压成几个数。压成几个数的时候，选择哪个数本身就是一种立场。

\subsection{频数表与图形}

\begin{definition}[频数表]
当数据只有少数几个不同取值时，把每个取值和它出现的次数（频数）列成表，就是频数表。
\end{definition}

\begin{example}[42 位毕业生的起薪]
起薪（千美元）与对应人数：
\[
\begin{array}{c|ccccccccccc}
x & 57 & 58 & 59 & 60 & 61 & 62 & 63 & 64 & 66 & 67 & 70\\ \hline
f & 4 & 1 & 3 & 5 & 8 & 10 & 0 & 5 & 2 & 3 & 1
\end{array}
\]
最低 57（4 人），最高 70（1 人），最常见的是 62（10 人）。一张表就把这批数据的形状交代完了。
\end{example}

\begin{definition}[相对频数]
大小为 $n$ 的数据集中，频数为 $f$ 的取值，其相对频数为 $f/n$。
\end{definition}

图形上有几种画法：条形图相当于加粗的折线图；频率多边形是把各柱顶用直线连起来；饼图用扇形\emph{面积}表示相对频数。

\begin{example}[200 位癌症患者的病种分布]
\[
\begin{array}{c|cc}
\text{类型} & \text{新发病例} & \text{相对频数}\\ \hline
\text{肺} & 42 & 0.21\\
\text{乳腺} & 50 & 0.25\\
\text{结肠} & 32 & 0.16\\
\text{前列腺} & 55 & 0.275\\
\text{黑色素瘤} & 9 & 0.045\\
\text{膀胱} & 12 & 0.06
\end{array}
\]
相对频数加起来是 $1$，所以饼图能画得圆满；这类构成性的问题，第一反应该是饼图而不是条形图。
\end{example}

\subsection{分组数据与直方图}

取值太多太散时，分组。

\begin{definition}[组距分组]
把取值范围切成若干个相邻区间（通常 5 到 10 组，长度多取相等），统计每组内的数据个数。区间写成左闭右开：$[500,600)$ 含 $500$、不含 $600$。
\end{definition}

\begin{example}[200 只白炽灯的寿命（组距 100）]
\[
\begin{array}{c|cccccc}
\text{区间} & [500,600) & [600,700) & [700,800) & [800,900) & [900,1000)\\ \hline
\text{频数} & 2 & 5 & 12 & 25 & 58
\end{array}
\qquad
\begin{array}{c|ccccc}
\text{区间} & [1000,1100) & [1100,1200) & [1200,1300) & [1300,1400) & [1400,1500)\\ \hline
\text{频数} & 41 & 43 & 7 & 6 & 1
\end{array}
\]
\end{example}

\begin{definition}[直方图]
以相邻摆放的矩形条表示各组频数的图形，就是直方图。它与条形图的区别正在「相邻」二字：条形图管的是离散类别，柱间有空隙；直方图管的是连续取值，柱子必须贴在一起。
\end{definition}

\begin{definition}[累计频数图]
把「小于某值的累计个数」逐点描出连成曲线，称为累计频数图（ogive）。它的用处是能直接从图上读出百分位数。
\end{definition}

\begin{definition}[茎叶图]
把每个数据拆成「茎」（高位）和「叶」（低位）两行写出的图。它的好处是既显示分布形状，又不丢掉任何原始数值——数据量中小时最划算。
\end{definition}

\begin{example}[35 个城市的年均最低气温（$^\circ$F）]
茎为十位，「$4\mid 3.6$」读作 $43.6$。
\[
\begin{array}{c|l}
\text{茎} & \text{叶}\\ \hline
7 & 0.0\\
6 & 9.0\\
5 & 1.0,\ 1.3,\ 2.0,\ 5.5,\ 7.1,\ 7.4,\ 7.6,\ 8.5,\ 9.3\\
4 & 0.0,\ 1.0,\ 2.4,\ 3.6,\ 3.7,\ 4.8,\ 5.0,\ 5.2,\ 6.0,\ 6.7,\ 8.1,\ 9.0,\ 9.2\\
3 & 3.1,\ 4.1,\ 5.3,\ 5.8,\ 6.2,\ 9.0,\ 9.5,\ 9.5\\
2 & 9.0,\ 9.8
\end{array}
\]
一眼能看出这批城市集中在 $40$ 多度，最长尾的是 $70$。
\end{example}

\subsection{集中趋势：均值、中位数、众数}

\begin{definition}[样本均值]
设样本为 $x_1,\dots,x_n$，
\[
\bar x=\frac1n\sum_{i=1}^{n}x_i .
\]
若数据来自频数表（$x_i$ 出现 $f_i$ 次），则 $\bar x=\frac1n\sum_i x_i f_i$。
\end{definition}

\begin{example}[先减掉一个基准数]
数据 $280,278,272,276,281,279,276,281,289,280$。逐个加太累，全部减去 $280$ 得
\[
0,-2,-8,-4,1,-1,-4,1,9,0,
\]
其平均为 $-8/10=-0.8$，所以 $\bar x=280-0.8=279.2$。手算时这一招能省下不少错误。
\end{example}

\begin{definition}[样本中位数]
把 $n$ 个值从小到大排好。若 $n=2k+1$，中位数是第 $k+1$ 个；若 $n=2k$，中位数是第 $k$ 与第 $k+1$ 个的平均。
\end{definition}

\begin{definition}[样本众数]
频数最大的那个取值称为众数；若几个取值并列最大，它们都是众数。
\end{definition}

\begin{example}[三个数一起算]
54 名学生的年龄分布：
\[
\begin{array}{c|cccccc}
\text{年龄} & 15 & 16 & 17 & 18 & 19 & 20\\ \hline
\text{频数} & 2 & 5 & 11 & 9 & 14 & 13
\end{array}
\]
均值
\[
\bar x=\frac{15\cdot2+16\cdot5+17\cdot11+18\cdot9+19\cdot14+20\cdot13}{54}=\frac{985}{54}\approx 18.24 .
\]
中位数是排序后第 27 与第 28 个的平均 $=18.5$。这题里两个数不一样，说明分布略向左偏（15 岁那一档比 20 岁那一档重）。
\end{example}

\begin{proposition}[均值与中位数的分工]
均值把每个数都用上了，因此对极端值敏感；中位数只看中间位置，极端值几乎不影响它。两者差得远，就是在提示你有尾巴。
\end{proposition}

\begin{example}[受体 300 rad 的小鼠]
比较无菌环境与常规环境下死于胸腺淋巴瘤的时间（天）：无菌组均值 $344.07$、中位数 $259$；常规组均值 $292.32$、中位数 $265$。均值差很多，中位数几乎一样。原因是无菌组里有 5 只活过了 500 天，把均值整条拉高，而中位数对这几只毫无反应。结论应当写成「无菌条件似乎只对活得最久的那部分个体有帮助」，而不是「平均多活 52 天」。
\end{example}

\begin{example}[掷骰子 40 次]
\[
\begin{array}{c|cccccc}
\text{点数} & 1 & 2 & 3 & 4 & 5 & 6\\ \hline
\text{频数} & 9 & 8 & 5 & 5 & 6 & 7
\end{array}
\]
众数 $=1$（出现 9 次）；均值 $=(9+16+15+20+30+42)/40=3.3$；中位数 $=3$（第 20、21 个的平均）。注意这里「众数 1」在骰子的语境里毫无意义——类别型的众数放到数值型数据上常常是这样，别当成中心位置。
\end{example}

\subsection{离散程度：样本方差与标准差}

描述了中心，还要描述散布。

\begin{definition}[样本方差与标准差]
\[
s^2=\frac{1}{n-1}\sum_{i=1}^{n}(x_i-\bar x)^2,\qquad s=\sqrt{s^2}.
\]
\end{definition}

\begin{proposition}[两个常用恒等式]
\[
\sum_{i=1}^{n}(x_i-\bar x)^2=\sum_{i=1}^{n}x_i^{2}-n\bar x^{2};
\]
若 $y_i=a+bx_i$，则 $\bar y=a+b\bar x$，$s_y^{2}=b^{2}s_x^{2}$。
\end{proposition}

第二个式子很有用：换单位（摄氏度对华氏、美元对千美元）时，均值跟着平移缩放，方差只被缩放因子的\emph{平方}放大，标准差正好被 $|b|$ 放大——所以想看原单位下的散布，要用标准差而不是方差。

\begin{remark}[为什么除以 $n-1$]
除以 $n-1$ 而不是 $n$，是为了让 $s^2$ 成为总体方差的无偏估计。这一点要等第七章讲到估计才说得清，现在先记住这个约定，并且知道「减 1」来自 $\bar x$ 本身已经吃掉了一个自由度。
\end{remark}

\begin{example}[同均值，不同散布]
$A:3,4,6,7,10$ 与 $B:-20,5,15,24$，两组均值都是 $6$。
\[
s_A^{2}=\frac{(-3)^2+(-2)^2+0^2+1^2+4^2}{4}=7.5,\qquad
s_B^{2}=\frac{(-26)^2+(-1)^2+9^2+18^2}{3}\approx 360.67 .
\]
光看均值完全分不出这两组，标准差立刻就分开了。这就是要造一个散布指标的原因。
\end{example}

\subsection{百分位数与四分位数}

\begin{definition}[样本 $100p$ 百分位数]
称某个值为样本 $100p$ 百分位数（$0\le p\le 1$），如果至少有 $100p\%$ 的数据不大于它，且至少有 $100(1-p)\%$ 的数据不小于它。若有两个数据值同时满足，就取两者的算术平均。
\end{definition}

实操规则（$n$ 为样本量）：
\begin{itemize}
  \item 若 $np$ 不是整数，取第 $\lceil np\rceil$ 个（向上取整的位置）；
  \item 若 $np$ 是整数，取第 $np$ 与第 $np+1$ 个的平均。
\end{itemize}

\begin{example}[美国人口最多的 25 个城市]
$n=25$，按人口从大到小排名 $1$ 到 $25$，则顺序统计量满足 $x_{(1)}\le\cdots\le x_{(25)}$，即 $x_{(k)}$ 对应排名第 $26-k$ 的城市。

第 10 百分位：$np=25\times0.10=2.5$ 不是整数，取第 $\lceil2.5\rceil=3$ 个（第 3 小，即排名第 23 的波士顿），$x_{(3)}=590{,}763$。

第 80 百分位：$np=25\times0.80=20$ 是整数，取第 20 与第 21 小的平均，即排名第 6 与第 5 的两个城市：
\[
\frac{1{,}448{,}394+1{,}512{,}986}{2}=1{,}480{,}690 .
\]
\end{example}

\begin{definition}[四分位数]
第 25 百分位记作 $Q_1$（第一四分位数），第 50 百分位就是中位数 $Q_2$，第 75 百分位记作 $Q_3$。
\end{definition}

\begin{definition}[箱线图]
把五数概括（最小值、$Q_1$、$Q_2$、$Q_3$、最大值）画在一条数轴上：盒子从 $Q_1$ 拉到 $Q_3$，中间用竖线标出中位数，两端的须伸到最小值与最大值。极差为 $\max-\min$，四分位距为 $\mathrm{IQR}=Q_3-Q_1$。
\end{definition}

\begin{example}[Grand Central 车站外的噪声]
36 次测量（dB）已排序。$np=36\times0.25=9$ 是整数，故
\[
Q_1=\frac{75+77}{2}=76,\quad
Q_2=\frac{89+90}{2}=89.5,\quad
Q_3=\frac{102+107}{2}=104.5 .
\]
极差 $=125-60=65$，$\mathrm{IQR}=104.5-76=28.5$。盒子里装着中间那 18 个测量值，须的长度则把两个极端交代干净。
\end{example}

\subsection{切比雪夫不等式}

前面那些数只说明「这一批数据」怎么样。切比雪夫给的是一个对任何数据都成立的保底比例。

\begin{theorem}[切比雪夫不等式]
设样本均值为 $\bar x$、样本标准差 $s>0$。对任意 $k\ge 1$，落在 $\bar x\pm ks$ 内的数据占比至少为
\[
100\left(1-\frac{1}{k^{2}}\right)\%.
\]
常用三档：$k=3/2$ 至少 $55.56\%$；$k=2$ 至少 $75\%$；$k=3$ 至少 $88.9\%$。
\end{theorem}

\begin{proof}[证明要点]
记 $S_k=\{i:|x_i-\bar x|<ks\}$，$|S_k|$ 为其元素个数。因为
\[
(n-1)s^{2}=\sum_{i=1}^{n}(x_i-\bar x)^{2}\ge\sum_{i\notin S_k}(x_i-\bar x)^{2}\ge k^{2}s^{2}\,(n-|S_k|),
\]
两边除以 $nk^{2}s^{2}$ 并整理，得
\[
\frac{|S_k|}{n}\ge 1-\frac{n-1}{nk^{2}}>1-\frac{1}{k^{2}} .
\]
\end{proof}

注意真正的界是 $1-\frac{n-1}{nk^{2}}$，比课本常写的 $1-1/k^{2}$ 略好；$n$ 大了两者几乎一样。另外 $k<1$ 时右边是负数，白说，所以定理只管 $k\ge1$。

\begin{proposition}[单边版本]
设 $N(k)=\#\{i:x_i-\bar x\ge ks\}$，$k\ge1$。由 $\{x_i-\bar x\ge ks\}\subseteq\{|x_i-\bar x|\ge ks\}$ 得
\[
\frac{N(k)}{n}\le\frac{1}{k^{2}} .
\]
更强的形式对任意 $k>0$ 成立：
\[
\frac{N(k)}{n}\le\frac{1}{1+k^{2}} .
\]
\end{proposition}

\begin{example}[2013 年 6 月畅销车型]
十款车的销量（千辆）给出 $\bar x=35.33$，$s=11.86$。取 $k=3/2$，切比雪夫保证区间
\[
\bar x\pm1.5s=(17.54,\ 53.12)
\]
内至少有 $55.56\%$ 的数据。实际上这批数据有 $90\%$ 落在里面。
\end{example}

这就是切比雪夫的性格：它什么分布都敢开口，所以给的界一定偏保守。想拿紧的界，就得对分布多作些假设——下一节就是最常用的一种假设。

\subsection{近似正态的数据集}

\begin{definition}[正态数据集与偏态]
直方图呈钟形、关于峰值对称的，称为正态型；接近这种形状的样本数据集称为近似正态。凡不关于中位数近似对称的，称为偏态：长尾在左叫左偏，长尾在右叫右偏。
\end{definition}

\begin{theorem}[经验法则]
若数据集近似正态，均值为 $\bar x$、标准差为 $s$，则
\begin{enumerate}
  \item 约 $68\%$ 的观测落在 $\bar x\pm s$；
  \item 约 $95\%$ 落在 $\bar x\pm2s$；
  \item 约 $99.7\%$ 落在 $\bar x\pm3s$。
\end{enumerate}
\end{theorem}

\begin{proposition}[两个工具的分工]
切比雪夫对任何数据都成立，但界松；经验法则给出的比例精确得多，前提是数据近似正态。判断能不能用经验法则，先看直方图或茎叶图有没有钟形和对称性，不要一拿到数就套 $68$--$95$--$99.7$。
\end{proposition}

\begin{example}[期末成绩]
某系工业工程学生的考试成绩（茎叶图，$n=28$），$\bar x\approx70.571$，$s\approx14.354$。
\[
\begin{array}{c|l}
\text{茎} & \text{叶}\\ \hline
9 & 0,\ 1,\ 4\\
8 & 3,\ 5,\ 5,\ 7,\ 8\\
7 & 2,\ 4,\ 4,\ 5,\ 7,\ 7,\ 8\\
6 & 0,\ 2,\ 3,\ 4,\ 6,\ 6\\
5 & 2,\ 5,\ 5,\ 6,\ 8\\
4 & 3,\ 6
\end{array}
\]
$\bar x\pm s\approx(56.2,84.9)$，$\bar x\pm2s\approx(41.9,99.3)$，$\bar x\pm3s\approx(27.5,113.6)$。数一下茎叶图里的落在各区间的个数，就能验证这三个比例是否接近 $68\%$、$95\%$、$99.7\%$。这份数据大致对称，验证下来相当服帖。
\end{example}

\subsection*{本章容易踩的地方}

\begin{itemize}
  \item 组区间是左闭右开，边界上的值只归一组；手抄数据时最容易在这一格数错。
  \item 方差除以 $n-1$，不是 $n$。混用会让后续所有涉及 $s^2$ 的结论差一个因子。
  \item 百分位数有好几套定义，讲义这套（看 $np$ 是否为整数）和统计软件默认的插值法不完全一致，对答案时要先统一口径。
  \item 均值和中位数不一致本身就是信息，别只报其中一个；有极端值时只报均值几乎一定误导。
  \item 类别型数据的众数不代表中心位置。
  \item 切比雪夫只管 $k\ge1$，而且它给的只是下界，实际比例通常高得多；经验法则反过来很准，但只对近似正态的数据成立。
\end{itemize}
